\chapter{Conclusiones y propuesta}
\label{cap:conclusiones}

\section{Conclusiones}

\begin{enumerate}
  \item La reproducción local de las mejores soluciones públicas de IEEE-CIS Fraud Detection es
        viable en hardware de consumidor (16 núcleos, RTX 2060 de 6\,GB) y entrega resultados
        cualitativamente idénticos a los de la competencia: el ordenamiento por AUC se conserva.
  \item El experimento de ablación demuestra que el rendimiento excepcional de la solución
        ganadora no proviene del algoritmo (idéntico en ambas variantes) ni del hardware, sino
        de la \textbf{ingeniería de variables orientada a entidades}: reconstruir el cliente
        implícito (\emph{magic UID}) y agregar su historial aporta $+0{,}0104$ de AUC local,
        más que cualquier cambio de modelo estudiado.
  \item La \textbf{estabilización temporal} de variables (normalización de columnas \texttt{D},
        filtrado por consistencia temporal) y la \textbf{validación temporal honesta}
        (\texttt{GroupKFold} mensual) son condiciones habilitantes: sin ellas, las ganancias no
        se transferirían al mes de prueba.
  \item El hardware acelera 4.4$\times$ el ciclo de experimentación sin mover el AUC: su valor es
        de productividad, no de precisión.
  \item Las técnicas de re-muestreo (SMOTE) presentes en el corpus analítico no forman parte de
        las soluciones ganadoras: con árboles y AUC, el desbalance se maneja vía objetivos
        probabilísticos y \texttt{scale\_pos\_weight}.
\end{enumerate}

\section{Propuesta de desarrollo}
\label{sec:futuro}

Sobre la evidencia recogida, la propuesta consiste en construir el pipeline siguiente, en orden
de retorno esperado:

\begin{enumerate}
  \item \textbf{Base}: XGBoost/LightGBM con las familias \texttt{C/D/V} seleccionadas por
        correlación y columnas \texttt{D} normalizadas (componentes ya verificados).
  \item \textbf{Entidades causales}: implementar el magic UID con agregaciones
        \emph{estrictamente históricas} (cada transacción usa solo estadísticas de
        transacciones anteriores a su \texttt{TransactionDT}), eliminando el carácter
        transductivo detectado en el protocolo original.
  \item \textbf{Ensemble}: combinar por promedio de rangos el XGB de entidades con un LightGBM
        equivalente ---el camino seguido por el equipo ganador (0{,}9459 en el tablero
        privado)--- y añadir el post-procesado de suavizado por UID.
  \item \textbf{Evaluación de negocio}: incorporar umbral de decisión y curva de utilidad
        (costo de bloqueo vs.\ pérdida por fraude) junto al AUC.
\end{enumerate}

Con ello, el proyecto transita de \emph{reproducir} las mejores técnicas públicas a
\emph{producir} una solución propia, superando la limitación metodológica principal detectada.
